\section{Introduction}

Formal-language-constrained path querying (FLPQ)~\cite{doi:10.1137/S0097539798337716} provides a mechanism for searching paths in edge-labeled graphs where constraints are specified in terms of formal languages.
A path is considered valid if the word formed by concatenating the labels of its edges belongs to the specified language.
When constraints are regular languages, such queries are called regular path queries (RPQs).
In practice, it is often useful to allow traversal opposite to edge directions, leading to two-way regular path queries (2-RPQs)~\cite{SIGMOD-Record-2003,DBLP:conf/amw/BienvenuCOS14}.
\cnote{C2}{Merge this and the next sentence into one (``Single-source and single-destination variants \ldots{} have been adopted in GQL, SQL/PGQ and SPARQL''). About 1 line.}
Queries specifying only a start or an end vertex are known as single-source or single-destination queries, respectively.
Both have been adopted in graph query languages such as GQL~\cite{2024gql}, SQL/PGQ, and SPARQL~\cite{2013sparql}.

\cnote{C3}{Cut this generic sentence and the last one of the paragraph (``Each semantics serves \ldots''); keep only the enumeration of semantics, and fix its all-shortest-paths gloss, which describes \textsc{any shortest}, not \textsc{all shortest}. About 3 lines.}
Graph database systems support different path query semantics, as the choice significantly affects both result interpretation and computational complexity.
Common semantics include: reachability (all vertices reachable via any valid path), all-shortest-paths (one shortest path per reachable vertex), all-simple-paths (all simple paths satisfying the query), and all-trails (all trails satisfying the query)~\cite{DBLP:conf/amw/FariasRV23}.
Each semantics serves different use cases, and the query engine must select the appropriate algorithm accordingly.

\cnote{C4}{Generic; fold into the next sentence (``Among the approaches to efficient (2-)RPQ evaluation~[..], sparse linear algebra has shown promise~[..]''). About 1 line.}
Efficient evaluation of RPQs and 2-RPQs remains an active research area \cite{casel_et_al:LIPIcs.ICDT.2021.19,DBLP:journals/corr/abs-2204-11137}.
Sparse linear algebra-based methods have shown promise for RPQ evaluation \cite{10.1007/978-3-031-43980-3_4}.
In particular, the LARPQ (Linear Algebra-based Regular Path Query) algorithm \cite{belyanin2150single} employs a breadth-first traversal that simultaneously processes both the graph and the finite automaton representing the regular language.
While LARPQ demonstrates promising performance, it was designed solely for reachability, and its application to path-related semantics (e.g., simple or shortest paths) remains unexplored.

The contributions of this paper are as follows.
\cnote{C5}{Contributions: bullet~2 can simply list the four instantiations; bullet~3 should be one line (the numbers belong to Sect.~6). About 3 lines. The introduction must also \emph{gain} a short delta-vs-[6] paragraph (ADBIS D7, R4) and ideally a two-sentence motivating example (R3), so expect the section to stay about the same length.}
\begin{itemize}
	\item A formal generalization of the LARPQ algorithm to a semiring-based framework, enabling the solution of various path-aggregation problems within a unified algorithmic structure.
	\item Instantiation of the framework to solve distinct problems: the Boolean semiring for reachability, and the path semiring with different index unary operators for the all-simple-paths, all-trails, and all-shortest-paths semantics.
	\item Experimental evaluation on both real-world and synthetic knowledge graphs and queries, showing that the approach is competitive with modern graph database systems: it is slower on the typical simple query but substantially faster on hard queries, which yields lower total and mean evaluation times.
\end{itemize}